\ssr{ВВЕДЕНИЕ}


В данной лабораторной работе рассматриваются графовые модели алгоритмов с целью выявления участков программного кода, которые могут быть выполнены параллельно. Анализ проводится на материале двух реализаций одной задачи --- рекурсивного и нерекурсивного алгоритмов, полученных из предыдущей лабораторной работы.

Для достижения поставленной цели необходимо решить следующие задачи:
\begin{enumerate}
	\item описать рекурсивную и нерекурсивную версии алгоритма согласно индивидуальному варианту;
	\item представить обе реализации с использованием четырёх графовых моделей --- графа управления, информационного графа, операционной истории и информационной истории;
	\item определить участки программного кода, которые могут быть выполнены параллельно, либо обосновать отсутствие таких участков;
\end{enumerate}




\clearpage


